\documentclass[12pt]{article}
\usepackage[utf8]{inputenc}
\usepackage[T1]{fontenc}
\usepackage[margin=1in]{geometry}
\usepackage{setspace}
\usepackage{xcolor}
\usepackage[colorlinks=true, urlcolor=blue!60!black, linkcolor=blue!60!black]{hyperref}
\setlength{\parindent}{2em}
\setlength{\parskip}{1.25em}
\renewcommand{\baselinestretch}{1.0}
\begin{document}
\begin{center}
{\LARGE \textit{Students for Fair Admissions, Inc. v. President and Fellows of Harvard College}, 600 U.S. 181 (2023)\par}
\vspace{0.8em}
{\large Case Brief --- Four-Page Edition\par}
\end{center}
\vspace{1em}

\section{Memory Jogger}

The Court invalidated the race-conscious undergraduate admissions programs at Harvard College and the University of North Carolina under the Equal Protection Clause and Title VI, holding that they lacked measurable objectives, used race as a negative and a stereotype, and had no logical endpoint, while allowing universities to consider an applicant's own discussion of how race affected their life (\href{https://www.law.cornell.edu/supremecourt/text/20-1199}{600 U.S.\ 181}).

\section{Detailed Case Facts}

Students for Fair Admissions (SFFA), a membership organization founded in 2014, sued Harvard, a private university subject to Title VI, and UNC, a public university subject to the Equal Protection Clause. At Harvard, first readers could consider race in an overall rating; committees reviewed ``one-pagers'' showing the racial makeup of the tentative class; and the final ``lop'' list showed only legacy status, athlete status, financial aid eligibility, and race. Black students made up 10.0\% to 11.7\% of admitted classes from 2009 to 2018, and the district court found that race was a ``determinative tip'' for many admitted Black and Hispanic applicants while resulting in fewer Asian American and White admissions. At UNC, readers could award a racial ``plus'' that ``may be significant in an individual case.'' Both schools relied on the diversity interest recognized in \textit{Bakke}, \textit{Grutter}, and \textit{Fisher}.

\section{Procedural History}

After a 15-day bench trial, the District of Massachusetts upheld Harvard's program, and the First Circuit affirmed. After an eight-day trial, the Middle District of North Carolina upheld UNC's. The Court granted certiorari in \textit{Harvard} and certiorari before judgment in \textit{UNC}, separating them so Justice Jackson could recuse from \textit{Harvard}. Argued October 31, 2022; decided June 29, 2023.

\section{Judicial Votes}

Reversed, \textbf{6--2} in \textit{Harvard} and \textbf{6--3} in \textit{UNC}.
\begin{itemize}
\item \textbf{Majority:} Chief Justice Roberts, joined by Justices Thomas, Alito, Gorsuch, Kavanaugh, and Barrett.
\item \textbf{Concurring:} Justice Thomas; Justice Gorsuch, joined by Justice Thomas; Justice Kavanaugh.
\item \textbf{Dissenting:} Justice Sotomayor, joined by Justice Kagan and, in \textit{UNC}, Justice Jackson; Justice Jackson in \textit{UNC}, joined by Justices Sotomayor and Kagan.
\end{itemize}

\section{Holding}

SFFA had associational standing under \textit{Hunt}. Both programs fail strict scrutiny: their diversity goals are too amorphous for judicial review, they use race as a ``negative'' in a zero-sum process, they rely on racial stereotyping, and they lack a logical endpoint. Universities may still consider an applicant's discussion of how race affected his or her life if the benefit is tied to that applicant's character or unique ability. The Court reserved military academies in footnote 4 and did not formally overrule \textit{Grutter}, though Thomas and Sotomayor agreed the decision ends race-based admissions in civilian higher education.

\section{Analysis of Opinions}

\textbf{Majority (Roberts).} Tracing equal protection from \textit{Strauder} and \textit{Yick Wo} to \textit{Brown}, Roberts held that the Clause is universal and that racial classifications survive only strict scrutiny, which outside remedying identified discrimination or preventing imminent prison violence the Court has almost never satisfied. Measured against \textit{Grutter}'s guardrails, the programs failed: goals like training leaders were ``commendable'' but unmeasurable, and the six racial categories were imprecise; because admissions are zero-sum, a plus for some is a minus for others, including Asian Americans; assuming race predicts viewpoint is stereotyping; and stable racial shares ``like clockwork'' showed no endpoint.

\textbf{Thomas (concurring).} Thomas defended a colorblind Constitution rooted in the Declaration, the 1866 Act, and original meaning, read the Freedmen's Bureau Acts as race-neutral, and called \textit{Grutter} effectively overruled.

\textbf{Gorsuch (concurring).} Title VI's text bars treating anyone worse ``on the ground of race,'' regardless of motive, as in \textit{Bostock}'s reading of Title VII.

\textbf{Kavanaugh (concurring).} The result follows from \textit{Grutter}'s own 25-year expectation; race-neutral outreach remains available.

\textbf{Sotomayor (dissenting).} The Court ``cement[s] a superficial rule of colorblindness'' and discards 45 years of precedent; the Reconstruction Congress enacted race-conscious Freedmen's Bureau Acts, and the trial courts found no quotas or anti-Asian discrimination.

\textbf{Jackson (dissenting in \textit{UNC}).} Jackson documented ``gulf-sized race-based gaps'' in wealth, health, and education traced to slavery, Jim Crow, and federal housing policy, and faulted the majority's ``let-them-eat-cake obliviousness.''

\section{Examples: Future Applications}

\begin{enumerate}
\item \textbf{Same-Side (Invalid): Medical School Target Bands.} A public medical school uses demographic dashboards and an informal plus to hold Black and Hispanic enrollment at 12\%--14\% with no sunset, replicating every defect condemned in \textit{SFFA}.
\item \textbf{Same-Side (Invalid): Race-Box Scholarship Points.} Automatic scholarship points for checking racial boxes violate Title VI, which tracks the Equal Protection Clause.
\item \textbf{Opposite-Side (Valid): Essay About Overcoming Discrimination.} Credit for an applicant's leadership in founding a legal-aid clinic after facing housing discrimination ties the benefit to her character, within the safe harbor.
\item \textbf{Opposite-Side (Valid): Class-Based Preferences.} A preference for low-income, first-generation applicants applied to all races is race-neutral and does not trigger strict scrutiny.
\item \textbf{Fence-Sitter: Military Academies.} West Point's use of race for officer diversity falls within footnote 4's reserved ``potentially distinct interests,'' though critics note most officers come from civilian ROTC programs now barred from using race.
\end{enumerate}

\section{Critique}

\textbf{Originalist and textualist views.} Originalists applaud the restoration of the anti-classification principle and Title VI's text, arguing that \textit{Grutter} invited racial balancing. Some note tensions: Roberts avoided original meaning, did not formally overrule \textit{Grutter}, and exempted military academies without explaining why colorblindness yields there.

\textbf{Progressive and doctrinal views.} Critics argue that race-conscious Reconstruction legislation refutes the colorblind reading and that the Fourteenth Amendment rests on an anti-subordination principle. They fault the Court for disregarding trial findings, unsettling 45 years of reliance, and inviting challenges to diversity programs and race-neutral reforms.

\section{Key Quotations}

\begin{quote}
``Eliminating racial discrimination means eliminating all of it. And the Equal Protection Clause, we have accordingly held, applies `without regard to any differences of race, of color, or of nationality'---it is `universal in [its] application.' ''
\end{quote}
\noindent --- Chief Justice Roberts, Opinion of the Court, 600 U.S.\ at 206.

\begin{quote}
``College admissions are zero-sum. A benefit provided to some applicants but not to others necessarily advantages the former group at the expense of the latter.''
\end{quote}
\noindent --- Chief Justice Roberts, Opinion of the Court, 600 U.S.\ at 218--219.

\begin{quote}
``At the same time, nothing prohibits universities from considering an applicant's discussion of how race affected the applicant's life, so long as that discussion is concretely tied to a quality of character or unique ability that the particular applicant can contribute to the university.''
\end{quote}
\noindent --- Chief Justice Roberts, Opinion of the Court, 600 U.S.\ at 230--231.

\end{document}
